\chapter{Imperium Inane}

\epigraf{Tots volen ser amos i cap l'amo de si mateix.}{Ugo Foscolo}

\lettrine[loversize=0.2, lines=2]{L}{a} paradoxa fonamental de la voluntat humana no es troba en el desig de llibertat, sinó en la seva substitució per la voluntat de domini. Els éssers humans no aspiren, al fons, a ser lliures: aspiren a governar. Però aquest govern que cerquen no és sobre si mateixos, sinó sobre allò que els envolta. I aquí rau el desajust: la conquesta del món ha esdevingut una fugida del govern intern. El desig de poder exterior no és sinó una expressió d’allò que no s’ha resolt dins nostre: una voluntat insatisfeta que no es reconeix en la pròpia limitació, sinó en el desig de superar allò que ens sembla inferior o feble.

Aquest desig de poder no és més que la projecció d’un buit interior, una recerca d’autoafirmació a través de la dominació de l’exterior. Els éssers humans han estat creats per lluitar contra els seus propis límits, per desitjar i aspirar més enllà del que poden ser. Però és a través de la confrontació amb les pròpies febleses, i no mitjançant la dominació d’allò que els envolta, que poden realitzar la seva veritable grandesa. Només quan un s’enfronta a la pròpia nul·litat, pot començar a crear-se com un ésser que s’autoafirma, que defineix la seva pròpia essència.

L’individu que vol imposar-se a l’altre sovint no és sinó aquell que ha fracassat en imposar-se a si mateix. El poder projectat cap a l'exterior actua com una ombra del buit interior. No es tracta d’una tirania conscient, sinó d’un desequilibri estructural: la manca de domini sobre la pròpia voluntat genera la necessitat de dominar l’entorn, com si aquest gest pogués compensar el que no s’ha resolt dins del propi ésser. La dominació de l’exterior només pot ser una resposta desesperada davant de la impotència interna. Qui no és capaç d’impo-\ sar-se sobre els propis desitjos, sobre les pròpies pors i dubtes, cerca refugi en l’autoritat que pot establir sobre els altres, buscant així la seguretat en un món de regles, jerarquies i estructures que li donen una sensació de control.

Però aquest desig de control no és només una necessitat personal, sinó també una estratègia de defensa davant de la por de l’incert. Si no podem controlar el nostre interior, almenys intentarem dominar allò que ens envolta. I així es crea una distorsió de l’ordre natural de les coses: aquells que no han conquerit la seva pròpia voluntat busquen conquerir el món extern, sense adonar-se que aquesta batalla és, en última instància, una guerra perduda. Perquè el control extern mai no pot substituir el control intern. El poder veritable neix d’aquell que ha estat capaç de transformar-se a si mateix, de dominar la pròpia natura, de superar les seves febleses i limitacions. La superació no és només un acte de força, sinó una contínua revalorització de les pròpies capacitats, una afirmació incessant de la vida mateixa.

El veritable senyor, però, no és aquell que imposa la seva voluntat als altres, sinó aquell que ha après a dominar les seves pròpies passions i desitjos. Aquell que s’ha disciplinat a si mateix no sent cap necessitat compulsiva d’exercir poder. Pot fer-ho, si cal, però no perquè en depengui. El veritable senyor és aquell que ha domesticat les seves pulsions, no pas qui imposa lleis a les dels altres. La llibertat veritable no és absència de dominació, sinó presència d’un ordre interior construït des del reconeixement i no des del rebuig. La llibertat es troba en el coratge de desafiar els propis límits i expandir la pròpia existència, no en l’intent de fugir d’ells.

Aquest capítol no defensa la submissió ni l’anul·lació de la voluntat. Defensa, en canvi, una inversió del desig: tornar l’atenció cap endins. No perquè l’interior sigui més pur, sinó perquè només en el silenci de la introspecció podem comprendre el veritable abast de la nostra llibertat. Fins que un no es posseeix a si mateix, tot poder sobre els altres serà, inevitablement, un reflex de la seva mancança. La veritable força no és la que es demostra en el control sobre el món, sinó la que es desplega en el control sobre l’individu mateix. Només en reconèixer les nostres pròpies debilitats, i després en superar-les, es pot obtenir un poder que no depèn de l’exterior, sinó de l’interior.

Potser la gran obra pendent de l’ésser humà no és l’organització externa del món, sinó la construcció sòbria del seu propi centre. I és només quan aquest centre esdevé estable que tot l’univers exterior pot ser comprès, no com un camp de batalla, sinó com un espai de correspondència. La vida no és un problema a resoldre, sinó un camp a conquerir. La recerca del poder no es basa en conquerir el món, sinó en conquerir-se un mateix, en realitzar les potencialitats que habiten dins de cada individu, en desfer-se de les antigues formes d’existència que no són sinó expressions d’una vida no viscuda plenament.

El domini veritable comença sempre per l’acceptació de la pròpia ignorància i per la difícil conquesta d’un mateix. Acceptar-se a un mateix és un acte revolucionari, que desafia les lleis establertes de la societat, que imposen una existència de conformisme i renúncia a la voluntat. La veritable llibertat es troba en l’afirmació de l’individu davant del món, en la capacitat de crear noves formes de vida i pensament, en no subordinat a les regles preexistents. Només quan un es reconeix a si mateix com a font de tot poder, pot conquerir el món, no com un conqueridor, sinó com un creador.